% Generated draft: requires mathematical and source review.
\documentclass[11pt]{article}
\usepackage[margin=25mm]{geometry}
\usepackage{fontspec}
\setmainfont{Latin Modern Roman}
\usepackage{amsmath,amssymb,mathtools,unicode-math}
\setmathfont{Latin Modern Math}
\usepackage{enumitem,needspace,xurl,hyperref}
\setlength{\parindent}{0pt}
\setlength{\parskip}{5pt}
\setlength{\emergencystretch}{3em}
\newcommand{\N}{\mathbb{N}}
\newcommand{\Z}{\mathbb{Z}}
\newcommand{\Q}{\mathbb{Q}}
\newcommand{\R}{\mathbb{R}}
\newcommand{\C}{\mathbb{C}}
\newcommand{\F}{\mathbb{F}}
\newcommand{\A}{\mathbb{A}}
\newcommand{\PP}{\mathbb{P}}
\newcommand{\E}{\mathbb{E}}
\newcommand{\eps}{\varepsilon}
\newcommand{\dd}{\,\mathrm{d}}
\begin{document}
\section*{UnsolvedMath ID: 6700042}
\textbf{Source problem code:} catalogue entry
\subsection{Definitions and mathematical statement}
For oriented smooth manifolds \(M,N\), a map \(f:M\to N\) is called a spin map, in the usual relative sense, when the stable oriented bundle \(TM\oplus f^*TN\) admits a spin structure. Its obstruction is
\[
 w_2(TM)=f^*w_2(TN)\quad\text{in }H^2(M;\Z/2\Z).
\]
This does not require \(M\) and \(N\) separately to be spin. For closed connected oriented manifolds of the same dimension, \(\deg f\) is defined by \(f_*[M]=(\deg f)[N]\).

Gromov's Spin Problem asks whether the spin-map hypothesis is necessary in the preceding scalar curvature extremality and rigidity statements, especially for equidimensional maps with \(\deg f\ne0\). The preceding statements are not included in this item, so their curvature assumptions and conclusions cannot be reconstructed as a single quantified theorem. The source explains that the hypothesis allows the required twisted Dirac operator to be defined; it does not assert that all nonzero-degree maps must be spin.

The appended passage makes a separate approximation assertion about ``Extremality via Lipschitz Maps'': if the indicated nonzero-degree, length decreasing, strictly scalar-curvature decreasing maps do not exist in class \(C^2\), then they do not exist in the Lipschitz class either. To distinguish the two metrics lost in extraction, write \(f:(M,g)\to(N,h)\). The usual nonexpanding interpretation is
\[
 d_h(f(x),f(y))\leq d_g(x,y),\qquad
 \operatorname{Scal}_h(f(x))<\operatorname{Scal}_g(x).
\]
For smooth maps the first condition is equivalent to \(f^*h\leq g\). These conventions, the source/target distinction, compactness assumptions, and whether ``length decreasing'' is strict require checking. The excerpt supplies no approximation details guaranteeing that its exact inequalities survive smoothing. The appended assertion is retained as a source report, not established here.
\subsection{Short English statement}
Can scalar curvature extremality and rigidity results for nonzero-degree maps dispense with their relative spin condition? The excerpt also reports a passage from smooth to Lipschitz maps by approximation, but omits the surrounding hypotheses.
\subsection{Sources}
\begin{itemize}
\item Local UnsolvedMath registry, permanent ID 6700042, source code catalogue entry; Misha Gromov, 2017; updated 2026-08-06.
\item Gromov, \emph{101 Questions, Problems and Conjectures around Scalar Curvature} (2017), printed tag [?45], occurrence 44, PDF page 52, ``Spin Problem'' and appended ``Extremality via Lipschitz Maps'': \url{https://www.ihes.fr/~gromov/wp-content/uploads/2018/08/101-problemsOct1-2017.pdf}. Registry access date 2026-07-29; normalized extraction from a public author-hosted incomplete PDF. Current status and reuse terms require review.
\item The supplied partial-progress summary mentions low-dimensional results, rational eta-invariants, and Dirac-index arguments without precise additional references or hypotheses. Its present-day status claims are not independently verified.
\end{itemize}
\end{document}
